\chapter{Software Architecture and Workflow}
\label{ch:software}

\section{The workflow end to end}

The project was built as a pipeline in which each stage is verified before the
next depends on it. Figure~\ref{fig:workflow} shows the whole chain.

\begin{figure}[htbp]
\centering
\begin{tikzpicture}[
  font=\small,
  node distance=6mm,
  stage/.style={draw, rounded corners=2pt, minimum height=9mm, text width=36mm, align=center},
  check/.style={draw, dashed, rounded corners=2pt, minimum height=8mm, text width=36mm, align=center, font=\footnotesize},
  >={Stealth[length=2.2mm]}
]
\node[stage] (ds) {Manufacturer datasheet\\ \footnotesize $\Voc,\Isc,\Vmp,\Imp,\beta_{Voc}$};
\node[stage, below=of ds] (fit) {Five-parameter identification\\ \footnotesize simultaneous nonlinear solve};
\node[stage, below=of fit] (pv) {PV array model\\ \footnotesize analytic $\mathrm{d}I/\mathrm{d}V$, exact MPP};
\node[stage, below=of pv] (avg) {Averaged converter model\\ \footnotesize state-space, linearised};
\node[stage, below=of avg] (ctrl) {Cascaded control design\\ \footnotesize margins over the envelope};
\node[stage, below=of ctrl] (mppt) {MPPT layer\\ \footnotesize P\&O, INC, ESC};
\node[stage, below=of mppt] (bench) {Dynamic benchmark\\ \footnotesize EN 50530-style ramps};
\node[stage, below=of bench] (sw) {Switched Simscape model\\ \footnotesize real devices at \SI{20}{\kilo\hertz}};
\node[stage, below=of sw] (hw) {Hardware plan\\ \footnotesize ratings, sensing, protection};

\foreach \a/\b in {ds/fit, fit/pv, pv/avg, avg/ctrl, ctrl/mppt, mppt/bench, bench/sw, sw/hw}
  \draw[->] (\a) -- (\b);

\node[check, right=14mm of fit] (c1) {residual at machine precision};
\node[check, right=14mm of pv]  (c2) {slope vs finite difference $<10^{-5}$};
\node[check, right=14mm of avg] (c3) {Jacobian vs finite difference $<10^{-9}$};
\node[check, right=14mm of ctrl](c4) {phase margin $>\SI{45}{\degree}$ everywhere};
\node[check, right=14mm of mppt](c5) {limit cycle $=2\Delta V$};
\node[check, right=14mm of sw]  (c6) {switched vs averaged $<\SI{0.25}{\percent}$};

\foreach \a/\b in {fit/c1, pv/c2, avg/c3, ctrl/c4, mppt/c5, sw/c6}
  \draw[->, dashed] (\a) -- (\b);
\end{tikzpicture}
\caption[Project workflow]{The project workflow. Solid arrows are the modelling chain; dashed boxes are the verification applied at each stage before the next one is allowed to depend on it.}
\label{fig:workflow}
\end{figure}

The organising principle is that \textbf{no stage is trusted until it is checked
against something independent}. The parameter fit is checked against the
datasheet; the analytic derivatives against finite differences; the averaged
model against a switched circuit simulation; the simplified design models
against the full-order plant. This matters because the errors that are dangerous
in a study like this are not crashes but plausible wrong numbers.

\section{Repository layout}

\begin{table}[htbp]
\centering
\caption{Repository structure.}
\label{tab:repo}
\small
\begin{tabular}{@{}>{\raggedright\arraybackslash}p{0.40\textwidth}>{\raggedright\arraybackslash}p{0.52\textwidth}@{}}
\toprule
\textbf{Path} & \textbf{Contents} \\
\midrule
\texttt{matlab/+mpptlab/+pv/} & \texttt{PVArray}, \texttt{extractParams}, \texttt{datasheet} \\
\texttt{matlab/+mpptlab/+conv/} & \texttt{Boost}: averaging, linearisation, envelope \\
\texttt{matlab/+mpptlab/+ctrl/} & \texttt{designCascade}, \texttt{loopAnalysis}, \texttt{PI}, \texttt{designSingleLoopPI} \\
\texttt{matlab/+mpptlab/+mppt/} & \texttt{Tracker}, \texttt{PandO}, \texttt{IncCond}, \texttt{ESC} \\
\texttt{matlab/+mpptlab/+bench/} & \texttt{simulate}, \texttt{irradiance}, \texttt{trackerSet} \\
\texttt{matlab/+mpptlab/+util/} & configuration, figure style, paths \\
\texttt{matlab/scripts/} & \texttt{s01} to \texttt{s05} analysis scripts \\
\texttt{matlab/tests/} & 27 unit tests in three classes \\
\texttt{matlab/run\_all.m} & runs the whole pipeline \\
\texttt{simulink/+pvsrc/pv\_module.ssc} & PV array as a Simscape component \\
\texttt{simulink/build\_mppt\_model.m} & builds the switched model from code \\
\texttt{simulink/run\_switched\_validation.m} & switched against averaged, closed loop \\
\texttt{results/figures/}, \texttt{results/data/} & generated outputs \\
\texttt{docs/report/} & this report \\
\bottomrule
\end{tabular}
\end{table}

\section{Design decisions worth justifying}

\subsection{A package namespace rather than loose functions}

All code lives under the \texttt{mpptlab} package. Calls read as
\texttt{mpptlab.pv.PVArray} and \texttt{mpptlab.ctrl.loopAnalysis}, which makes
provenance obvious at the call site and prevents name collisions with MATLAB
functions such as \texttt{simulate} and \texttt{config}.

\subsection{Handle classes for physical objects}

\texttt{PVArray} and \texttt{Boost} are handle classes. A PV array is a physical
thing with mutable state, its irradiance and temperature, and passing it by
reference into a simulation loop that updates those properties thousands of
times per second is both natural and efficient.

\subsection{One configuration function}

Every script obtains its array, converter and gains from
\texttt{mpptlab.util.config}. There is exactly one place where the module type,
the array topology, the component values and the crossover frequencies are
defined. Changing the array from six modules to eight is a one-line edit that
propagates correctly through every figure and table.

\subsection{A common tracker interface}

All trackers derive from an abstract \texttt{Tracker} class exposing
\texttt{reset} and \texttt{update}. The base class owns reference clamping and
update counting; subclasses implement only \texttt{advance}. This guarantees the
comparison is fair: every tracker sees the same plant, the same inner loops, the
same sample rate mechanics and the same limits, so measured differences are
attributable to the algorithm and nothing else.

\subsection{The Simulink model is generated by a script}

\texttt{build\_mppt\_model.m} constructs the model block by block and wires it
programmatically. A \texttt{.slx} file is a binary archive: it cannot be
meaningfully diffed, reviewed or merged. A builder script can. The model is
therefore reproducible from source, and a change to a component value or a
control gain appears in version control as a readable one-line change.

\section{Numerical methods}

\subsection{Solving the diode equation quickly}

The implicit diode equation must be solved several times per integration step.
Three measures make the closed-loop simulations tractable: warm starting from
the previous solution, which cuts Newton to two or three iterations; clamping
the exponent argument to avoid overflow during transients; and a backtracking
line search that guarantees the residual magnitude never grows.

\subsection{Integration}

The averaged plant is integrated with fixed-step fourth-order Runge--Kutta at
half the switching period, \SI{25}{\micro\second}, with the discrete controllers
updating every \SI{50}{\micro\second}. A step-size convergence test is part of
the test suite: halving the step changes the final voltage by under
\SI{0.5}{\volt}, confirming the result is not an artefact of the integrator.

\subsection{Scoring}

The available power used in the efficiency denominator comes from the exact MPP
solver evaluated on an irradiance grid and interpolated, as described in
Chapter~\ref{ch:benchmark}.

\section{Verification}

\subsection{The test suite}

Twenty-seven tests run in three classes.

\begin{table}[htbp]
\centering
\caption{Test coverage.}
\label{tab:tests}
\begin{tabular}{@{}p{0.30\textwidth}rp{0.48\textwidth}@{}}
\toprule
\textbf{Class} & \textbf{Tests} & \textbf{What is checked} \\
\midrule
\texttt{tPVModel} & 11 & datasheet reproduction, short and open circuit conditions, MPP stationarity, residual on the whole curve, analytic slope against finite difference, monotonicity, irradiance and temperature trends, array scaling, $r_d = \Vmp/\Imp$ at the MPP, zero-irradiance robustness \\
\texttt{tConverterAndControl} & 9 & operating point is an equilibrium, Jacobian against finite differences, power gain vanishes at the MPP and changes sign across it, modal formulas against eigenvalues, sign of the duty-to-voltage gain, cascade margins above \SI{45}{\degree} over the envelope, single-loop margin below \SI{10}{\degree}, continuous conduction mode \\
\texttt{tMPPTAndBench} & 7 & profile shapes and slopes, closed-loop reference holding against the analytic I--V curve, limit cycle equals twice the step, all trackers converge to the MPP, efficiency bounded by one, integrator step convergence \\
\bottomrule
\end{tabular}
\end{table}

Several of these are not routine. The test that the cascade holds above
\SI{45}{\degree} while the single-loop design falls below \SI{10}{\degree}
encodes the central engineering conclusion of the project as an executable
assertion, so a future change that quietly destroys the argument fails a test.
Likewise the test that the duty-to-power gain vanishes at the MPP pins the
theoretical result of Section~\ref{sec:vanishing}.

\subsection{Reproducing everything}

\begin{table}[htbp]
\centering
\caption{Result of the full end-to-end run.}
\label{tab:runall}
\begin{tabular}{@{}lrr@{}}
\toprule
\textbf{Stage} & \textbf{Status} & \textbf{Time (\si{\second})} \\
\midrule
s01 PV characterisation   & OK   & 14.5 \\
s02 plant analysis        & OK   & 5.8 \\
s03 control design        & OK   & 9.0 \\
s04 MPPT comparison       & OK   & 39.9 \\
s05 dynamic benchmark     & OK   & 356.4 \\
s06 switched validation   & OK   & 156.9 \\
s07 model diagram export  & OK   & 6.3 \\
Unit tests                & 27 / 27 passed & 42.0 \\
\midrule
\textbf{Total}            &      & \textbf{630.8} \\
\bottomrule
\end{tabular}
\end{table}

The command \texttt{run\_all} executes the entire pipeline, regenerates every
figure and table in this report, rebuilds the Simulink model from its builder
script, and runs the test suite. Table~\ref{tab:runall} is the result of that
run on the development machine. Every number quoted anywhere in this report
comes from that pipeline; none were transcribed by hand from an earlier
experiment.
